\documentclass[12pt]{article}
\usepackage{amsmath,amssymb,amsthm}
\usepackage[margin=1in]{geometry}
\usepackage{algorithm,algpseudocode}

\newtheorem{theorem}{Theorem}
\newtheorem{lemma}[theorem]{Lemma}
\newtheorem{definition}[theorem]{Definition}

\title{Problem 182: RSA Encryption}
\author{Project Euler Solution}
\date{}

\begin{document}
\maketitle

\section{Problem Statement}

Given distinct primes $p = 1009$ and $q = 3643$, let $n = pq$ and $\lambda(n) = \operatorname{lcm}(p-1, q-1)$. Find the sum of all valid RSA exponents $e$ (with $1 < e < \lambda(n)$, $\gcd(e, \lambda(n)) = 1$) that minimise the number of unconcealed messages $m \in \{0, \ldots, n-1\}$ satisfying $m^e \equiv m \pmod{n}$.

\section{Mathematical Development}

\begin{theorem}[Unconcealed Message Count]
For RSA modulus $n = pq$ with distinct odd primes $p, q$ and exponent $e$ with $\gcd(e, \lambda(n)) = 1$:
\[
N(e) = \bigl(1 + \gcd(e-1, p-1)\bigr)\bigl(1 + \gcd(e-1, q-1)\bigr).
\]
\end{theorem}

\begin{proof}
By the Chinese Remainder Theorem, $m^e \equiv m \pmod{n}$ iff $m^e \equiv m \pmod{p}$ and $m^e \equiv m \pmod{q}$. For a prime $r \in \{p, q\}$, the equation $m(m^{e-1} - 1) \equiv 0 \pmod{r}$ has two disjoint solution sets: $m \equiv 0 \pmod{r}$ (one solution) and $m^{e-1} \equiv 1$ in $(\mathbb{Z}/r\mathbb{Z})^*$. By Lagrange's theorem applied to the cyclic group of order $r - 1$, the latter has exactly $\gcd(e-1, r-1)$ solutions. The total modulo $r$ is $1 + \gcd(e-1, r-1)$, and by CRT the answer is the product.
\end{proof}

\begin{theorem}[Minimisation]
The minimum of $N(e)$ over valid exponents is $9$, achieved when $\gcd(e-1, p-1) = 2$ and $\gcd(e-1, q-1) = 2$.
\end{theorem}

\begin{proof}
Since $\lambda(n)$ is even, every valid $e$ is odd, so $e - 1$ is even, forcing $\gcd(e-1, p-1) \ge 2$ and $\gcd(e-1, q-1) \ge 2$. Setting $a = \gcd(e-1, p-1)$ and $b = \gcd(e-1, q-1)$, the function $(1+a)(1+b)$ with $a, b \ge 2$ is minimised at $a = b = 2$, giving $9$.

Existence: with $p-1 = 2^4 \cdot 3^2 \cdot 7$ and $q-1 = 2 \cdot 3 \cdot 607$, the conditions $v_2(e-1) = 1$, $3 \nmid (e-1)$, $7 \nmid (e-1)$, $607 \nmid (e-1)$, and $\gcd(e, \lambda(n)) = 1$ are simultaneously satisfiable by CRT (the moduli $4, 3, 7, 607$ are pairwise coprime).
\end{proof}

\begin{lemma}[Explicit Conditions]
An exponent $e$ achieves $N(e) = 9$ if and only if: $\gcd(e, \lambda(n)) = 1$, $e \equiv 3 \pmod{4}$, $e \not\equiv 1 \pmod{3}$, $e \not\equiv 1 \pmod{7}$, and $e \not\equiv 1 \pmod{607}$.
\end{lemma}

\begin{proof}
These conditions are equivalent to $\gcd(e-1, p-1) = 2$ and $\gcd(e-1, q-1) = 2$ after factoring $p - 1 = 2^4 \cdot 3^2 \cdot 7$ and $q - 1 = 2 \cdot 3 \cdot 607$, as verified prime-by-prime.
\end{proof}

\section{Algorithm}

\begin{algorithmic}[1]
\Function{SumOptimalExponents}{$p, q$}
    \State $\lambda_n \gets \operatorname{lcm}(p-1, q-1)$;\quad $\mathrm{total} \gets 0$
    \For{$e = 2$ to $\lambda_n - 1$}
        \If{$\gcd(e, \lambda_n) \neq 1$} \textbf{continue} \EndIf
        \If{$\gcd(e-1, p-1) = 2$ \textbf{and} $\gcd(e-1, q-1) = 2$}
            \State $\mathrm{total} \mathrel{+}= e$
        \EndIf
    \EndFor
    \State \Return total
\EndFunction
\end{algorithmic}

\section{Complexity Analysis}

\begin{itemize}
    \item \textbf{Time:} $O(\lambda(n) \log \lambda(n))$, iterating over all $e < \lambda(n) = 612{,}648$ with $O(\log \lambda(n))$ GCD cost per candidate.
    \item \textbf{Space:} $O(1)$.
\end{itemize}

\section{Answer}

\[
\boxed{399788195976}
\]

\end{document}
